\section*{\texorpdfstring{\S\CurrentExerciseGroup}{第{\CurrentExerciseGroup}节}}
\phantomsection
\addcontentsline{toc}{section}{第{\CurrentExerciseGroup}节习题}

¶ 1) a) 设 X 是一个局部紧空间。试证：局部凸空间 \(\mathcal{H}(\mathrm{X};\mathbf{C})\) 的每个包含于 \(\mathcal{H}_{+}(\mathrm{X})\) 中的紧凸子集 A 都是严格紧的。（考虑所有 \(f\in\mathrm{A}\) 所成之集的上包络函数 \(s=\sup_{f\in\mathrm{A}}f\) ，它在

X 上连续且在 X 的无穷远点处趋于 0。对 X 上每个连续数值函数 g，用 \(\mathrm{P}(g)\) 表示由满足下述条件的 \(x \in X\) 所成之集：\(g(x) > 0\) 。试证 A 中存在一个函数 f 使 \(\mathrm{P}(f) = \mathrm{P}(s)\) ；为此，注意若 \(K_{n}\) 是由满足下述条件的 \(x \in X\) 组成的紧集：\(s(x) \geqslant 1/n\)，则存在有限个函数 \(f_{n,i} \in A\) 使诸开集 \(\mathrm{P}(f_{n,i})\) 覆盖 \(K_{n}\) 。然后利用 A 是凸的且闭的这一事实。）

\begin{enumerate}
\def\labelenumi{\alph{enumi})}
\setcounter{enumi}{1}
\tightlist
\item
  假设 X 中存在一个可数稠密子集 D。试证 \(\mathcal{K}(X; \mathbf{C})\) 的每个紧凸子集 A 此时都是严格紧的。（设 U 是由满足下述条件的 \(x \in X\) 所成之集：对某个 \(f \in A\) 有 \(f(x) \neq 0\) ，它是 X 中的开集。利用 A 是 Baire 空间这一事实，试证存在函数 \(g \in A\) 使 \(g(x) \neq 0\) 于所有 \(x \in U \cap D\) 处成立。）
\end{enumerate}

\begin{enumerate}
\def\labelenumi{\arabic{enumi})}
\setcounter{enumi}{1}
\tightlist
\item
  设 Y 是一个非紧且无穷远处可数的局部紧空间，并设 \((Y_{n})\) 是 Y 的紧子集的递增序列，构成 Y 的覆盖且 \(Y_{n} \subset Y_{n+1}\) 对所有 n 成立，使得 Y 的每个紧子集都含于某个 \(Y_{n}\) 之中（GT, I, §9, No.~9, 命题 15 的推论 1）。设 \(\beta Y\) 是 Y 的 Stone--Čech 紧化（GT, IX, §1, 习题 7），则空间 \(\mathcal{C}(\beta Y; \mathbf{R})\) 可等同于 Y 上有界连续实值函数所成的空间，且 Y 是 \(\beta Y\) 中的稠密开集；对每个函数 \(f \in \mathcal{C}(\beta Y; \mathbf{R})\) ，\(\operatorname{Supp}(f)\) 就是 \(\beta Y\) 中由满足下述条件的 \(y \in Y\) 组成之集的闭包：\(f(y) \neq 0\) 。回顾一下，若 \(Z_{1}, Z_{2}\) 是 Y 的两个在 Y 中闭且不相交的子集，则闭包 \(\overline{Z}_{1}, \overline{Z}_{2}\) 在 \(\beta Y\) 中也不相交（GT, IX, §4, 习题 17）。
\end{enumerate}

设 \(\omega\) 是 \(\beta\mathrm{Y}-\mathrm{Y}\) 的一个点，并令 \(\mathrm{X}=\beta\mathrm{Y}-\{\omega\}\) 。试证：在空间 \(\mathcal{K}(\mathrm{X};\mathbf{R})\) 上，No.~1 中定义的直接极限拓扑 \(\mathcal{T}\) 与一致收敛拓扑相同。（若 \(p\) 是对 \(\mathcal{T}\) 连续的一个半范数，用反证法证明：存在数 \(c > 0\) 使 \(p(f) \leqslant c \| f \|\) 对每个函数 \(f \in \mathcal{K}(\mathrm{X};\mathbf{R})\) 成立。否则，将存在函数的一个序列 \((f_n)\) （诸函数均属于 \(\mathcal{K}(\mathrm{X};\mathbf{R})\) ） ，使 \(\| f_n \| \leqslant 1\) ，\(p(f_n) \geqslant n\) ，\(\operatorname{Supp}(f_n) \cap Y_n = \emptyset\) ，且

\[
\operatorname{Supp} (f _ {n}) \cap \operatorname{Supp} (f _ {k}) = \varnothing
\]

对 \(k < n\) 成立。然后设 \(Z_{i}\) ( \(i = 0,1\) ) 是诸集 \(Y \cap \operatorname{Supp}(f_{2n + i})\) 的并；试证两集 \(Z_{0}, Z_{1}\) 之一在 X 中相对紧（GT, X, §4, 习题 7），并由此导出矛盾。）

由此推出：空间 \(\mathcal{K}(\mathbf{X};\mathbf{R})\) 赋以拓扑 \(\mathcal{T}\) 后不是拟完备的。

\begin{enumerate}
\def\labelenumi{\arabic{enumi})}
\setcounter{enumi}{2}
\tightlist
\item
  在习题 2 的例子中，取 Y 为离散空间 N。对每个 \(n \in N\) ，设 \(e_n\) 是 \(\mathcal{K}(X; \mathbf{R})\) 中满足 \(e_n(n) = 1\) 且 \(e_n(x) = 0\) （对 X 中 \(x \neq n\) ）的函数。试证由函数 0 和 \(e_n/(n + 1)\) 组成的集 A 在 \(\mathcal{K}(X; \mathbf{R})\) 中是紧的，但不是严格紧的。若 U 是由 \(\omega\) 在 \(\beta Y = \beta N\) 中的邻域滤子所诱导的 N 上的超滤子，试证：关于 U，正测度族 ( \(n\varepsilon_n\) ) 在 \(\mathcal{K}(X; \mathbf{R})\) 的每个严格紧子集上一致趋于 0，但在 A 中不一致趋于 0。
\end{enumerate}

¶ 4) a) 设 \(\mathbf{B}_n\) 是欧氏空间 \(\mathbf{R}^n\) 中的闭单位球，并设 \(\mathbf{A}_n\) 是这样的函数 \(f \in \mathcal{C}(\mathbf{B}_n; \mathbf{R})\) 所成的集：它形如

\[
f (x) = \lambda (\| x \| ^ {2} - 1) + (x | a)
\]

（其中 \(|\lambda| \leqslant 1/n\) 且 \(\|a\| \leqslant 1/n\) ）。集 \(A_n\) 在 Banach 空间 \(\mathcal{C}(\mathbf{B}_n; \mathbf{R})\) 中是凸的且紧的，并具有下述性质：对任意函数 \(f_1, \ldots, f_n\) 属于 \(A_n\) 者，存在点 \(x \in B_n\) 使 \(f_1(x) = \ldots = f_n(x) = 0\) ；然而，存在 \(n + 1\) 个函数 \(g_1, \ldots, g_{n+1}\) 属于 \(A_n\) ，它们在

\(\mathbf{B}_n\) 的任何点处都不同时为零（可以考虑映射 \(x \mapsto x / (1 - \|x\|^2)\) ，它把 \(\mathbf{B}_n\) 映入 \(\mathbf{R}^n\) ）。 b) 设 \(\mathbf{B}_n'\) 是赋以离散拓扑的集 \(\mathbf{B}_n\) ，并令 \(\mathrm{D}_n = \beta \mathrm{B}_n'\) 。设 Y （相应地 \(\mathrm{Y}'\) ）是诸空间 \(\mathrm{D}_n\) （相应地 \(\mathrm{B}_n'\) ）的局部紧拓扑和空间。试证 \(\beta \mathrm{Y} = \beta \mathrm{Y}'\) 。

\begin{enumerate}
\def\labelenumi{\alph{enumi})}
\setcounter{enumi}{2}
\tightlist
\item
  对每个整数 \(n\) ，设 \(\mathrm{K}_n \subset \mathcal{C}(\mathrm{D}_n; \mathbf{R})\) 是这样的实值函数全体所成的集：它们在 \(\mathrm{B}_n' = \mathbf{B}_n\) 上的限制属于 \(\mathrm{A}_n\) 。设 \(\mathrm{K}\) 是 \(\mathcal{C}(\beta\mathrm{Y}; \mathbf{R})\) 的由满足以下条件的函数组成的子集：其在 \(\mathrm{D}_n\) 上的限制属于 \(\mathrm{K}_n\) 对每个 \(n \in \mathbf{N}\) 成立。试证 \(\mathrm{K}\) 是 Banach 空间 \(\mathcal{C}(\beta\mathrm{Y}; \mathbf{R})\) 的一个紧凸子集。对函数的每个有限族 \((f_i)_{1 \leqslant i \leqslant n}\) （诸函数属于 \(\mathrm{K}\) ），设 \(W(f_1, \ldots, f_n)\) 是由满足下述条件的 \(y \in Y'\) 所成之集：\(f_i(y) = 0\) （对 \(1 \leqslant i \leqslant n\) ）。试证诸 \(W(f_1, \ldots, f_n)\) 构成某个滤子 \(\mathfrak{W}\) 在离散空间 \(Y'\) 上的基。由此得出：存在点 \(\omega \in \beta\mathrm{Y} - \mathrm{Y}\) ，使得在 \(Y'\) 上，由 \(\omega\) 在 \(\beta\mathrm{Y}\) 中的邻域滤子诱导的滤子是比 \(\mathfrak{W}\) 更细的滤子。令 \(X = \beta Y - \{\omega\}\) ，由此推出 \(K \subset \mathcal{K}(X; \mathbf{R})\) ，并借助习题 2 断定：在 \(\mathcal{K}(X; \mathbf{R})\) 中，集 \(K\) 是凸的且紧的，但不是严格紧的。
\end{enumerate}

\begin{enumerate}
\def\labelenumi{\arabic{enumi})}
\setcounter{enumi}{4}
\tightlist
\item
  试证：若 \(\mathbf{E}\) 是 Fréchet 空间且 \(\mathbf{X}\) 是局部紧空间，则空间 \(\mathcal{H}(\mathbf{X};\mathbf{E})\) 是有界型的。
\end{enumerate}

由此推出：对 \(\mathcal{K}(X;C)\) 的有界子集组成的每个集 S ，若它覆盖 \(\mathcal{K}(X;C)\) ，且 \(\mathcal{K}(X;C)\) 中由一个收敛序列的点所成的每个集都属于 S ，则空间 \(\mathcal{M}(X_{\mathfrak{S}};C)\) 是完备的。

\begin{enumerate}
\def\labelenumi{\arabic{enumi})}
\setcounter{enumi}{5}
\tightlist
\item
  设 \(\mu\) 是局部紧空间 \(X\) 上的一个测度。试证 \(|\mu| = \sup \mathcal{R}(\alpha \mu)\) ，其中标量 \(\alpha\) 遍历绝对值 \(\leqslant 1\) 的复数所成之集。（注意，对每个函数 \(g \in \mathcal{K}(X; \mathbf{C})\) 及每个 \(\varepsilon > 0\) ，存在有限个映射 \(h_i\) （把 \(X\) 映入 \([0,1]\) ，具有紧支集）以及标量 \(\alpha_i\) ，使 \(\| g - \sum_{i} \alpha_i h_i \| \leqslant \varepsilon\) 且 \(\| |g| - \sum_{i} |\alpha_i| h_i \| \leqslant \varepsilon\) 。）
\end{enumerate}

试证亦有 \(|\mu| = \sup \mathcal{R}(h \cdot \mu)\) ，其中 \(h\) 遍历 \(\mathcal{K}(\mathrm{X};\mathbf{C})\) 中满足 \(\| h\| \leqslant 1\) 的函数所成之集。

\begin{enumerate}
\def\labelenumi{\arabic{enumi})}
\setcounter{enumi}{6}
\item
  举一个例子说明：在 \(\mathcal{M}(\mathrm{X};\mathbf{C})\) 中，有界实测度 \(\mu\) 中满足 \(\|\mu\|=a\) 者所成之集未必是淡闭的，即使 X 是紧的也是如此；试证：若 X 是局部紧且非紧的，则满足 \(\|\mu\|=1\) 的有界正测度所成之集不是淡闭的。
\item
  设 \(\mathbf{X}\) 是一个局部紧空间；对每个紧子集 \(\mathbf{K}\) ⊂ \(\mathbf{X}\) 及每个整数 \(n > 0\) ，设 \(S_{\mathbf{K},n}\) 是由满足下述条件的函数 \(f \in \mathcal{K}(\mathbf{X};\mathbf{C})\) 组成的集：支集含于 \(\mathbf{K}\) 且 \(\| f\| \leqslant n\) 。\(\mathcal{M}(\mathbf{X};\mathbf{C})\) 上的拟强拓扑指的是在诸集 \(S_{\mathbf{K},n}\) 上一致收敛的拓扑。它比强拓扑粗（当 \(\mathcal{M}(\mathbf{X};\mathbf{C})\) 被视为空间 \(\mathcal{K}(\mathbf{X};\mathbf{C})\) 的对偶时，参见 TVS, III, §3, No.~1），而当 \(\mathbf{X}\) 仿紧时二者相同。
\end{enumerate}

\begin{enumerate}
\def\labelenumi{\alph{enumi})}
\item
  试证：\(\mathcal{M}(\mathrm{X};\mathbf{C})\) 上的拟强拓扑由诸半范数 \(\mu \mapsto |\mu |(f)\) 定义，其中 \(f\) 遍历 \(\mathcal{K}_{+}(\mathrm{X})\) ；空间 \(\mathcal{M}(\mathrm{X};\mathbf{C})\) 对于拟强拓扑是完备的。\(\mathcal{M}(\mathrm{X};\mathbf{C})\) 中关于拟强拓扑的有界子集就是 \(\mathcal{M}(\mathrm{X};\mathbf{C})\) 中的等度连续子集。
\item
  试证：在 \(\mathcal{M}(\mathrm{X};\mathbf{R})\) 上，拟强拓扑与序向量空间结构相容（TVS, II, §2, No.~7）。由此推出：若 H 是 \(\mathcal{M}(\mathrm{X};\mathbf{R})\) 中有上界的递增有向集，则 H 在 \(\mathcal{M}(\mathrm{X};\mathbf{R})\) 中的上确界与其截面滤子关于拟强拓扑的极限相同（loc. cit.）。
\end{enumerate}

\begin{enumerate}
\def\labelenumi{\arabic{enumi})}
\setcounter{enumi}{8}
\tightlist
\item
  设 \(\mathbf{X}\) 是紧区间 {[}0,1{]} ，它含于 \(\mathbf{R}\) 。
\end{enumerate}

\begin{enumerate}
\def\labelenumi{\alph{enumi})}
\item
  设 \(\mu_{n}\) 是由点 0 处的质量 \(+1\) 与质量 \(-1\) （置于点 \(1/n\) 处）所定义的测度。试证序列 \((\mu_{n})\) 在 \(\mathcal{M}(\mathrm{X};\mathbf{C})\) 中淡收敛于 0 ，但 \(|\mu_n|\) 不淡收敛于 0 。
\item
  设 \(\mu\) 是 X 上的 Lebesgue 测度，并令 \(g_{n}(x) = \sin nx\) 。试证测度 \((1 - g_{n}) \cdot \mu\) 是正的，且淡收敛于 \(\mu\) （当 \(n\) 趋于无穷时），但不强收敛于 \(\mu\) 。由此推出：淡拓扑与强拓扑（这里它与拟强拓扑相同）在 \(\mathcal{M}_{+}(X)\) 上诱导的拓扑是不同的（参见习题 10）。
\end{enumerate}

\begin{enumerate}
\def\labelenumi{\arabic{enumi})}
\setcounter{enumi}{9}
\tightlist
\item
  设 \(\Phi\) 是集 \(\mathcal{M}_{+}(\mathrm{X})\) （即 \(\mathrm{X}\) 上正测度全体所成之集）上的一个滤子。试证：若 \(\Phi\) 淡收敛，则对每个紧子集 \(\mathrm{K}\) ⊂ \(\mathrm{X}\) ，存在一个集 \(\mathrm{M} \in \Phi\) 及一个数 \(a_{\mathrm{K}} > 0\) ，使 \(|\mu(f)| \leqslant a_{\mathrm{K}} \cdot \|f\|\) 对每个函数 \(f \in \mathcal{K}(\mathrm{X}, \mathrm{K}; \mathbf{C})\) 及每个测度 \(\mu \in \mathrm{M}\) 成立。
\end{enumerate}

¶ 11) a) 试证：当 \(\mathcal{M}(\mathrm{X};\mathbf{C})\) 赋以拟强拓扑（相应地，严格紧收敛拓扑；相应地，淡拓扑）时，双线性型 \((f,\mu)\mapsto\langle f,\mu\rangle\) 是 \((\mathfrak{S},\mathfrak{T})\) -亚连续的，其中 \(\mathfrak{T}\) 是 \(\mathcal{M}(\mathrm{X};\mathbf{C})\) 的淡有界子集所成之集，而 \(\mathfrak{S}\) 是 \(\mathcal{H}(\mathrm{X};\mathbf{C})\) 中含于某个 \(\mathcal{H}(\mathrm{X},\mathrm{K};\mathbf{C})\) （K 为适当的紧集）的有界子集所成之集（相应地，\(\mathcal{H}(\mathrm{X};\mathbf{C})\) 的严格紧子集所成之集；相应地，\(\mathcal{H}(\mathrm{X};\mathbf{C})\) 的有限子集所成之集）（参见 TVS, III, §5, 习题 12）。

\begin{enumerate}
\def\labelenumi{\alph{enumi})}
\setcounter{enumi}{1}
\item
  对每个紧子集 \(\mathbf{K}\) ⊂ \(\mathbf{X}\) ，试证：双线性型 \((f, \mu) \mapsto \langle f, \mu \rangle\) 在 \(\mathcal{H}(\mathbf{X}, \mathbf{K}; \mathbf{C}) \times \mathcal{M}(\mathbf{X}; \mathbf{C})\) 上连续（当 \(\mathcal{M}(\mathbf{X}; \mathbf{C})\) 赋以拟强拓扑时）；它在 \(\mathcal{H}(\mathbf{X}, \mathbf{K}; \mathbf{C}) \times \mathcal{M}_{+}(\mathbf{X})\) 上连续（当 \(\mathcal{M}_{+}(\mathbf{X})\) 赋以淡拓扑时）（利用习题 10）。
\item
  假设 \(X\) 是仿紧且非紧的。试证：双线性型 \((f, \mu) \mapsto \langle f, \mu \rangle\) 在 \(\mathcal{K}(X; \mathbf{C}) \times \mathcal{M}_{+}(X)\) 上不连续（当 \(\mathcal{M}_{+}(X)\) 赋以强拓扑时，此处强拓扑与拟强拓扑相同）。
\item
  举一个例子：X 是紧的，且双线性型 \((f,\mu)\mapsto\langle f,\mu\rangle\) 在 \(\mathcal{C}(\mathrm{X};\mathbf{C})\times\mathcal{M}(\mathrm{X};\mathbf{C})\) 上不连续（当 \(\mathcal{M}(\mathrm{X};\mathbf{C})\) 赋以紧收敛拓扑时）（TVS, IV, §1, 习题 11）。
\end{enumerate}

¶ 12) a) 试证：双线性映射 \((g, \mu) \mapsto g \cdot \mu\) （从 \(\mathcal{C}(\mathrm{X}; \mathbf{C}) \times \mathcal{M}(\mathrm{X}; \mathbf{C})\) 到 \(\mathcal{M}(\mathrm{X}; \mathbf{C})\) ）当 \(\mathcal{C}(\mathrm{X}; \mathbf{C})\) 赋以紧收敛拓扑且 \(\mathcal{M}(\mathrm{X};\mathbf{C})\) 赋以拟强拓扑（习题 8）或严格紧收敛拓扑时是连续的。

\begin{enumerate}
\def\labelenumi{\alph{enumi})}
\setcounter{enumi}{1}
\item
  试证：双线性映射 \((g,\mu)\mapsto g\cdot \mu\) （从 \(\mathcal{C}(\mathrm{X};\mathbf{C})\times \mathcal{M}(\mathrm{X};\mathbf{C})\) 到 \(\mathcal{M}(\mathrm{X};\mathbf{C})\) ）是 \((\mathfrak{S},\mathfrak{T})\) -亚连续的（当 \(\mathcal{C}(\mathrm{X};\mathbf{C})\) 赋以紧收敛拓扑且 \(\mathcal{M}(\mathrm{X};\mathbf{C})\) 赋以淡拓扑时），其中 \(\mathfrak{S}\) 表示 \(\mathcal{C}(\mathrm{X};\mathbf{C})\) 的有限子集所成之集，\(\mathfrak{T}\) 表示 \(\mathcal{M}(\mathrm{X};\mathbf{C})\) 的淡有界子集所成之集。
\item
  试证：映射 \((g, \mu) \mapsto g \cdot \mu\) （从 \(\mathcal{C}(\mathrm{X};\mathbf{C}) \times \mathcal{M}_{+}(\mathrm{X})\) 到 \(\mathcal{M}(\mathrm{X};\mathbf{C})\) ）当 \(\mathcal{C}(\mathrm{X};\mathbf{C})\) 赋以紧收敛拓扑、\(\mathcal{M}(\mathrm{X};\mathbf{C})\) 与 \(\mathcal{M}_{+}(\mathrm{X})\) 赋以淡拓扑时是连续的。
\item
  试给出一个紧空间 X 的例子，使得双线性映射 \((g,\mu)\mapsto g\cdot\mu\) （从 \(\mathcal{C}(\mathrm{X};\mathbf{C})\times\mathcal{M}(\mathrm{X};\mathbf{C})\) 到 \(\mathcal{M}(\mathrm{X};\mathbf{C})\) ）当 \(\mathcal{C}(\mathrm{X};\mathbf{C})\) 赋以一致收敛拓扑且 \(\mathcal{M}(\mathrm{X};\mathbf{C})\) 赋以淡拓扑时不连续。
\end{enumerate}

\begin{enumerate}
\def\labelenumi{\arabic{enumi})}
\setcounter{enumi}{12}
\tightlist
\item
  试证：对每个函数 \(f \in \mathcal{K}_{+}(X)\) ，映射 \(\mu \mapsto |\mu|(f)\) 在 \(\mathcal{M}(X; \mathbf{C})\) 上关于淡拓扑是下半连续的。
\end{enumerate}

¶ 14) a) 设 X 是具有可数基的局部紧空间。试证：空间 \(\mathcal{M}_{+}(X)\) 赋以淡拓扑后是可度量的。

\begin{enumerate}
\def\labelenumi{\alph{enumi})}
\setcounter{enumi}{1}
\tightlist
\item
  设 L 是由这样的递增序列 \((k_{n})_{n\in\mathbf{Z}}\) 组成的集：诸项为整数 \(\geqslant0\) ，随 n 趋于 \(+\infty\) ，且除有限多个 n 的值外，在 \(n\leqslant0\) 时为零；设 P 是由平移 \((k_{n+h})_{n\in\mathbf{Z}}\) （取自同一个序列 \((k_{n})_{n\in\mathbf{Z}}\) ，其中 h 遍历 Z）所组成的等价类所成之集；L 与 P 都是不可数的。设 X 是 P 与 Z 的离散拓扑和空间；设 H 是 X 上按下述方式定义的正测度所成之集：对每个 \(\alpha\in P\) 及每个序列 \(g\in\alpha\) ，\(\nu_{\alpha,g}\) 是由点 \(\alpha\) 处的质量 +1 、P 的其他点处的质量 0 以及质量 \(g(n)\) （置于每个点 \(n\in\mathbf{Z}\) 处）所定义的测度；取 H 为由诸测度 \(\nu_{\alpha,g}\) 生成的锥。现在设 \(\mu\) 是 X 上由 P 的每个点处的质量 +1 与 Z 的每个点处的质量 0 所定义的测度。试证 \(\mu\) 属于 H 的淡闭包，但不属于 H 的任何淡有界子集的淡闭包。（利用下述事实：对 Z 到 N 中的每个映射 f ，存在递增序列 \((k_{n})\in L\) ，使对每个 \(h\in Z\) 有 \(\lim_{n\to\infty}k_{n+h}/f(n)=+\infty\) 。）
\end{enumerate}

\begin{enumerate}
\def\labelenumi{\arabic{enumi})}
\setcounter{enumi}{14}
\tightlist
\item
  设 \(\mathbf{X}\) 是非紧的局部紧空间；在有界测度所成的空间 \(\mathcal{M}^1(\mathbf{X};\mathbf{C})\) （即 \(\mathbf{X}\) 上关于 Banach 空间 \(\mathcal{C}^0(\mathbf{X};\mathbf{C})\) 的对偶）上，由范数 \(\| \mu \|\) 定义的拓扑称为超强拓扑，而拓扑 \(\sigma (\mathcal{M}^1(\mathbf{X};\mathbf{C}),\mathcal{C}^0(\mathbf{X};\mathbf{C}))\) 称为弱拓扑。
\end{enumerate}

\begin{enumerate}
\def\labelenumi{\alph{enumi})}
\item
  试证：若 \(X\) 仿紧，则 \(\mathcal{M}^1(X; \mathbf{C})\) 上的弱拓扑严格细于淡拓扑（与习题 2 比较）。
\item
  试证：在 \(\mathcal{M}^1 (\overline{\mathrm{X}};\mathbf{C})\) 上，超强拓扑严格细于拟强拓扑（注意，对后者而言，\(\varepsilon_{x}\) 当 \(x\) 趋于 X 的无穷远点时趋于 0 ）。
\item
  若 \(\mathbf{X}\) 在无穷远处可数且非离散，试证：在 \(\mathcal{M}^1(\mathbf{X};\mathbf{C})\) 上，弱拓扑与拟强拓扑不可比较。
\item
  在 \(\mathcal{M}^1 (\mathrm{X};\mathbf{C})\) 中，每个弱有界集关于超强拓扑有界，但可能存在这样的集：它关于拟强拓扑有界而非弱有界。
\end{enumerate}

\begin{enumerate}
\def\labelenumi{\arabic{enumi})}
\setcounter{enumi}{15}
\item
  设 X 是一个局部紧空间，\(X'\) 是在 X 上添加一个无穷远点所得到的紧空间。设 \(\mu\) 是 X 上的一个有界测度；试证：存在一个且仅一个测度 \(\mu'\) 于 \(X'\) 上，它延拓 \(\mu\) 且满足 \(\|\mu'\| = \|\mu\|\) 。
\item
  对每个测度 \(\mu\) （在局部紧空间 \(X\) 上）及每个同胚 \(\sigma\) （\(X\) 到其自身的），设 \(\mu_{\sigma}\) 为测度 \(f \mapsto \mu(f \circ \sigma)\) 。
\end{enumerate}

\begin{enumerate}
\def\labelenumi{\alph{enumi})}
\item
  试证 \(|\mu_{\sigma}| = |\mu|_{\sigma}\) 。
\item
  给 \(\mathcal{M}(\mathrm{X};\mathbf{C})\) 赋以淡拓扑，并用 A 表示局部凸空间 \(\mathcal{M}(\mathrm{X};\mathbf{C})\) 的连续自同态所成之集；给 A 赋以在 \(\mathcal{M}(\mathrm{X};\mathbf{C})\) 的淡有界子集上一致收敛的拓扑。设 \(\Gamma\) 是 X 到其自身的同胚所成的群，赋以 GT, X, §3, No.~5 中定义的拓扑 \(\mathcal{T}_{\beta}\) 。对每个 \(\sigma \in \Gamma\) ，设 \(A_{\sigma}\) 为映射 \(\mu \mapsto \mu_{\sigma}\) ，它属于 \(\mathcal{A}\) ；试证映射 \(\sigma \mapsto A_{\sigma}\) （从 \(\Gamma\) 到 \(\mathcal{A}\) 中）是连续的。
\end{enumerate}

\begin{enumerate}
\def\labelenumi{\arabic{enumi})}
\setcounter{enumi}{17}
\tightlist
\item
  设 \(X\) 是紧空间，\(\mu\) 是 \(X\) 上满足 \(\mu(1) = \| \mu \|\) 的测度。试证 \(\mu\) 是正测度。（若 \(\mu_1 = \mathcal{R}(\mu)\) ，\(\mu_2 = \mathcal{I}(\mu)\) ，先注意 \(\mu_1(1) = \| \mu_1 \|\) ，并利用 Ch. II, §2, No.~2, 命题 4 证明 \(\mu_1\) 是正测度；然后证明 \(\mu_2 = 0\) ，办法是考虑 \(\mu(1 + if)\) （对适当的函数 \(f \in \mathcal{C}(X; \mathbf{R})\) 。）
\end{enumerate}

\setcounter{exercisegroup}{2}
\edef\CurrentExerciseGroup{\arabic{exercisegroup}}
